\subsection{Camada de validação}

A calibração do validador (Seção~\ref{subsubsec:calibracao}) foi conduzida sobre um conjunto de \textit{anotações-ouro} pequeno e anotado por uma única pessoa, montado em separado para cada funcionalidade. Esse recorte é suficiente para demonstrar o método --- mostrar que a confiança verbalizada pode ser testada quanto à sua discriminância e usada para calibrar um limiar de abstenção ---, mas não para fixar valores de produção: tanto o limiar $\tau$ quanto as métricas que o sustentam (AUC e a curva de risco--cobertura) são ilustrativos, e não estimativas estáveis do desempenho do validador. Um conjunto maior, com vários anotadores e cobertura mais ampla de casos difíceis, seria necessário para reduzir a variância dessas medidas e para estabelecer a concordância entre anotadores como teto realista de desempenho --- referência contra a qual o acerto do validador deveria ser lido, já que em afirmações genuinamente ambíguas nem mesmo avaliadores humanos concordam plenamente entre si.

Uma segunda limitação é de escopo de implementação. A verificação de afirmações de conhecimento de mundo --- fatos históricos, lugares reais e pano de fundo cultural que vivem fora da obra --- foi especificada no desenho do validador, mas não implementada no presente trabalho, por exigir uma fonte externa (busca na web) com desafios próprios que extrapolam o TRL 3 visado (Seção~\ref{subsubsec:arquitetura}). Na prática, toda afirmação desse tipo é tratada conservadoramente como \textit{Não verificável} e conduz a resposta ao estado \textit{Com ressalvas}. A escolha é deliberadamente segura --- prefere-se admitir a incapacidade de confirmar a arriscar uma validação infundada ---, mas tem um custo: respostas factualmente corretas sobre o mundo real nunca chegam ao estado \textit{Válido}, o que pode subestimar a confiabilidade do assistente justamente nas contextualizações, em que esse tipo de afirmação é mais frequente.